%%%%%%%%%%%%%%%%%%%%%%%%%%%%%%%%%%%%%%%%%
% Programming/Coding Assignment
% LaTeX Template
%
% This template has been downloaded from:
% http://www.latextemplates.com
%
% Original author:
% Ted Pavlic (http://www.tedpavlic.com)
%
% This template uses a Perl script as an example snippet of code, most other
% languages are also usable. Configure them in the "CODE INCLUSION 
% CONFIGURATION" section.
%
%%%%%%%%%%%%%%%%%%%%%%%%%%%%%%%%%%%%%%%%%

%----------------------------------------------------------------------------------------
%	PACKAGES AND OTHER DOCUMENT CONFIGURATIONS
%----------------------------------------------------------------------------------------

\documentclass{article}
%\documentclass[english,a4paper,twoside]{amsart}

\usepackage{fancyhdr} % Required for custom headers
\usepackage{lastpage} % Required to determine the last page for the footer
\usepackage{extramarks} % Required for headers and footers
\usepackage[usenames,dvipsnames]{color} % Required for custom colors
\usepackage{graphicx} % Required to insert images
\usepackage{listings} % Required for insertion of code
\usepackage{courier} % Required for the courier font

\usepackage[utf8]{inputenc}
\usepackage[T1]{fontenc}

% Margins
\topmargin=-0.45in
\evensidemargin=0in
\oddsidemargin=0in
\textwidth=6.5in
\textheight=9.0in
\headsep=0.25in

\linespread{1.1} % Line spacing

% Set up the header and footer
\pagestyle{fancy}
\lhead{\hmwkAuthorName} % Top left header
\chead{\hmwkClass : \hmwkTitle} % Top center head
\rhead{\firstxmark} % Top right header
\lfoot{\lastxmark} % Bottom left footer
\cfoot{} % Bottom center footer
\rfoot{Page\ \thepage\ of\ \protect\pageref{LastPage}} % Bottom right footer
\renewcommand\headrulewidth{0.4pt} % Size of the header rule
\renewcommand\footrulewidth{0.4pt} % Size of the footer rule

\setlength\parindent{0pt} % Removes all indentation from paragraphs

%----------------------------------------------------------------------------------------
%	CODE INCLUSION CONFIGURATION
%----------------------------------------------------------------------------------------

\definecolor{MyDarkGreen}{rgb}{0.0,0.4,0.0} % This is the color used for comments
\lstloadlanguages{Perl} % Load Perl syntax for listings, for a list of other languages supported see: ftp://ftp.tex.ac.uk/tex-archive/macros/latex/contrib/listings/listings.pdf
\lstset{%language=Perl, % Use Perl in this example
        frame=single, % Single frame around code
        basicstyle=\small\ttfamily, % Use small true type font
        keywordstyle=[1]\color{Blue}\bf, % Perl functions bold and blue
        keywordstyle=[2]\color{Purple}, % Perl function arguments purple
        keywordstyle=[3]\color{Blue}\underbar, % Custom functions underlined and blue
        identifierstyle=, % Nothing special about identifiers                                         
        commentstyle=\usefont{T1}{pcr}{m}{sl}\color{MyDarkGreen}\small, % Comments small dark green courier font
        stringstyle=\color{Purple}, % Strings are purple
        showstringspaces=false, % Don't put marks in string spaces
        tabsize=5, % 5 spaces per tab
        %
        % Put standard Perl functions not included in the default language here
        morekeywords={rand},
        %
        % Put Perl function parameters here
        morekeywords=[2]{on, off, interp},
        %
        % Put user defined functions here
        morekeywords=[3]{test},
       	%
        morecomment=[l][\color{Blue}]{...}, % Line continuation (...) like blue comment
        numbers=left, % Line numbers on left
        firstnumber=1, % Line numbers start with line 1
        numberstyle=\tiny\color{Blue}, % Line numbers are blue and small
        stepnumber=5 % Line numbers go in steps of 5
}

% Creates a new command to include a perl script, the first parameter is the filename of the script (without .pl), the second parameter is the caption
\newcommand{\perlscript}[2]{
\begin{itemize}
\item[]\lstinputlisting[caption=#2,label=#1]{#1}
\end{itemize}
}

%----------------------------------------------------------------------------------------
%	DOCUMENT STRUCTURE COMMANDS
%	Skip this unless you know what you're doing
%----------------------------------------------------------------------------------------

% Header and footer for when a page split occurs within a problem environment
\newcommand{\enterProblemHeader}[1]{
\nobreak\extramarks{#1}{#1 continued on next page\ldots}\nobreak
\nobreak\extramarks{#1 (continued)}{#1 continued on next page\ldots}\nobreak
}

% Header and footer for when a page split occurs between problem environments
\newcommand{\exitProblemHeader}[1]{
\nobreak\extramarks{#1 (continued)}{#1 continued on next page\ldots}\nobreak
\nobreak\extramarks{#1}{}\nobreak
}

\setcounter{secnumdepth}{0} % Removes default section numbers
\newcounter{homeworkProblemCounter} % Creates a counter to keep track of the number of problems

\newcommand{\homeworkProblemName}{}
\newenvironment{homeworkProblem}[1][Problem \arabic{homeworkProblemCounter}]{ % Makes a new environment called homeworkProblem which takes 1 argument (custom name) but the default is "Problem #"
\stepcounter{homeworkProblemCounter} % Increase counter for number of problems
\renewcommand{\homeworkProblemName}{#1} % Assign \homeworkProblemName the name of the problem
\section{\homeworkProblemName} % Make a section in the document with the custom problem count
\enterProblemHeader{\homeworkProblemName} % Header and footer within the environment
}{
\exitProblemHeader{\homeworkProblemName} % Header and footer after the environment
}

\newcommand{\problemAnswer}[1]{ % Defines the problem answer command with the content as the only argument
\noindent\framebox[\columnwidth][c]{\begin{minipage}{0.98\columnwidth}#1\end{minipage}} % Makes the box around the problem answer and puts the content inside
}

\newcommand{\homeworkSectionName}{}
\newenvironment{homeworkSection}[1]{ % New environment for sections within homework problems, takes 1 argument - the name of the section
\renewcommand{\homeworkSectionName}{#1} % Assign \homeworkSectionName to the name of the section from the environment argument
\subsection{\homeworkSectionName} % Make a subsection with the custom name of the subsection
\enterProblemHeader{\homeworkProblemName\ [\homeworkSectionName]} % Header and footer within the environment
}{
\enterProblemHeader{\homeworkProblemName} % Header and footer after the environment
}

%----------------------------------------------------------------------------------------
%	NAME AND CLASS SECTION
%----------------------------------------------------------------------------------------

\newcommand{\hmwkTitle}{Miniproject\ \#1} % Assignment title
\newcommand{\hmwkClass}{Modeling and Verification} % Course/class
\newcommand{\hmwkClassInstructor}{Anna Ingolfsdóttir} % Teacher/lecturer
\newcommand{\hmwkAuthorName}{Þröstur and Sævar} % Your name

%----------------------------------------------------------------------------------------
%	TITLE PAGE
%----------------------------------------------------------------------------------------

\title{
\vspace{2in}
\textmd{\textbf{\hmwkClass:\ \hmwkTitle}}\\
\vspace{0.1in}\large{\textit{\hmwkClassInstructor}}
\vspace{3in}
}

\author{\textbf{\hmwkAuthorName}}
\date{} % Insert date here if you want it to appear below your name

%----------------------------------------------------------------------------------------

\begin{document}

\maketitle

%----------------------------------------------------------------------------------------
%	TABLE OF CONTENTS
%----------------------------------------------------------------------------------------

%\setcounter{tocdepth}{1} % Uncomment this line if you don't want subsections listed in the ToC

\newpage

%----------------------------------------------------------------------------------------
%	PROBLEM 1
%----------------------------------------------------------------------------------------

% To have just one problem per page, simply put a \clearpage after each problem

\begin{homeworkProblem}

  Task 1 was to implement a scheduler in ccs\ldots

Listing \ref{../task1.ccs} shows the CCS code for the scheduler.

\perlscript{../task1.ccs}{Scheduler .CCS file}

The scheduler was verifier with the file shown in listing \ref{../task1.mu} and foo bar bloo blah\ldots

\perlscript{../task1.mu}{Scheduler .MU file}

\end{homeworkProblem}

\newpage

%----------------------------------------------------------------------------------------
%	PROBLEM 2
%----------------------------------------------------------------------------------------

\begin{homeworkProblem}

  The second task at hand was to implement one-dimensional solitaire and to prove that it is solvable using the CCS concurrency workbench.
  The game was modelled as processes, where each tile is a process \(B(lack), W(hite),\) or \( E(mpty) \) and the board is a set of tiles (concurrent processes). 
  The three different colors were modelled as different processes. The processes were defined as templates and were created with appropriate relabellings to indicate location-based pairing between processes. 
	Solitaire has the helpful property that the $Empty$ tile must be swapped with whatever tile wants to swap, so we modelled the $Empty$ tile as having output ports for each direction and colored tiles to have input ports in their respective directions 
	($Black$ tiles can only go right and $White$ tiles can only go left) for the adjacent tile and its immediate neighbor.
	\emph{(Note: inputs and outputs can be freely interchanged here, since it is only the pairing that matters \-- $Empty$ tiles were selected as the output to improve readability in the source code.)}

	Listing \ref{../task2.ccs} demonstrates this functionality:

\perlscript{../task2.ccs}{1-Dimensional Solitaire .CCS file}
  Each template behaves differently based on what color it represents. Black tiles can jump either one or two tiles to the right depending on whether the $Empty$ tile is one or two tiles to the right (white tiles mirror this functionality in the opposite direction). Recall that $Empty$ tiles can swap with any adjacent tile.
  Relabelling the processes ensured that changing a process (for example, \(B \rightarrow W\)) allowed new input/output pairs to be matched.


  To transition from the configuration \([B_1|B_2|E_3|B_4|W_5|W_6|W_7]\) to \([B_1|B_2|W_3|B_4|E_5|W_6|W_7]\), $W_5$ must switch places with $E_3$ (by jumping over $B_4$).
  This move is made possible by the consistent location-dependent relabelling \-- recall that the board is modelled as concurrent processes.
  The relabelling on $E_3$ was \([b3/b, r4/rone, r5/rtwo, l3/l, r3/r, w3/w, l2/lone, l1/lone]\) when it was created as $B_3$, which means that $E_3$ exists as:
\[
	proc ~ E_3 = ~ \overline{r3}.B_3 + ~ \overline{l3}.W_3
\]
In a similar fashion, $W_5$ exists as:
\[
  proc ~ W_5 = ~ w_5.W_5 + ~ l4.E_5 + ~ l3.E_5
\]
Since these processes are running in parallel and restrictions are in place, $E_3$ is blocked until it is matched on either $'r3'$ or $'l3$. Incidentally, $W_5$ can match with $l3$, allowing the processes to swap 'colors' (the only other possibility would be for $W_5$ to indicate that tile 5 is white with the $w_5$ message).

In a similar fashion, if we had started in the configuration \([B_1|B_2|E_3|W_4|B_5|W_6|W_7]\), $W_4$ would be able to swap places with $E_3$ on $l3$ since $W_4$ would have been:
\[
	proc ~= W_4 = ~ w_4.W_4 + ~ l3.E_4 + ~ l2.E_4
\]

% CORRECTNESS AND VERIFICATION DETAILS BELOW

  Showing possible moves does not prove that the game is solvable (even though we can use these moves to solve it in our heads).
  To verify the correctness of the solitaire game, we simply asked if a configuration representing a solved game was reachable: 
  
\perlscript{../task2.mu}{1-Dimensional Solitaire .MU file}

  This elegant solution works due to the way that the game was modelled.
	Colored tiles (\(B_i, W_i\)) can always indicate what color occupies tile $i$ since $proc ~ Board$ does not restrict on $W_i$ (for any $i$).
	Therefore we can describe the property of a $solved$ game as ``Tiles $1-3$ are white and $5-7$ are black''.
	Using the fixed point-notation of Hennessy–Milner logic, we simply ask if the state is reachable from the start configuration:
	\[
		\$~cwb-nc> ~ chk~Board~is\_solvable
	\]

	\ldots which the concurrency workbench indicates as being $TRUE$.

% \problemAnswer{
% \begin{center}
% % \includegraphics[width=0.75\columnwidth]{example_figure} % Example image
% \end{center}
% 
% }
\end{homeworkProblem}

%----------------------------------------------------------------------------------------

\end{document}

